\documentclass[11pt]{article}
\usepackage{amsmath,amssymb,amsthm,hyperref}
\newtheorem{theorem}{Theorem}
\newtheorem{proposition}[theorem]{Proposition}
\newtheorem{remark}[theorem]{Remark}
\newcommand{\Q}{\mathbb Q}
\begin{document}
\title{An everywhere locally soluble one-coordinate pole curve}
\author{}\date{17 November 2026}\maketitle
Call $d\in\Q$ prepared if
\[
 0<d<1,\quad \operatorname{den}(d)\text{ is a square},\quad
 d\in7\Q_2^{\times2},\quad v_2(d)\ge4.
\]
Consider the genus-one curve
\begin{equation}\label{eq:C}
 C_d:\quad y^2=d\bigl((1-d)^2-s^4\bigr).
\end{equation}

\begin{theorem}\label{thm:local}
For every prepared input $d$, the smooth projective model of $C_d$ has a point over every completion of $\Q$. At each completion, points can be chosen in the affine chart with $sy\ne0$.
\end{theorem}
\begin{proof}
Write $d=A/B^2$ in lowest terms, with $A,B$ positive integers, and set $C=B^2-A$. Then $A,C>0$, $\gcd(A,C)=1$, $B$ is odd, $16\mid A$, and $C\equiv1\pmod{16}$. The change
\[
 t=Bs,\qquad Y=B^3y
\]
gives
\begin{equation}\label{eq:integral}
 Y^2=A(C^2-t^4).
\end{equation}
Its defining quartic has four distinct roots in characteristic zero, so its smooth projective model has genus one.

For an odd prime $p\mid A$, the identity $C=B^2-A$ implies that $C$ is a square in $\Q_p$. There is a local point $t=\sqrt C,Y=0$. For an odd prime $p\mid C$, $A=B^2-C$ is a square in $\Q_p$, and $t=0,Y=C\sqrt A$ is a local point. If $p\nmid2AC$, the quartic has smooth reduction of genus one over $\mathbb F_p$. The Hasse bound gives at least $p+1-2\sqrt p>0$ points on that reduction, and a smooth point lifts to $\Q_p$.

At $2$, the number $C\equiv1\pmod{16}$ is a square, so the first branch-point construction applies. At the real place take $t=0,Y=C\sqrt A$. Every point used is smooth on the curve over the relevant local field. A smooth local curve has infinitely many points in a sufficiently small analytic neighbourhood, so the finite set $tY=0$ and the points at infinity can all be avoided.
\end{proof}

\begin{proposition}\label{prop:sound}
Let $d\in\Q\setminus\{0,1\}$. If \eqref{eq:C} has a rational affine point, then $v_p(d)$ is never negative and odd, for any prime $p$.
\end{proposition}
\begin{proof}
Suppose $n=v_p(d)<0$ is odd. Then $v_p(1-d)=n$. If $s\ne0$, the two terms $d(1-d)^2$ and $-ds^4$ on the right have valuations $3n$ and $n+4v_p(s)$. Both are odd and they cannot be equal, since equality would require $n=2v_p(s)$. Their sum therefore has odd valuation, contradicting that it is a square. If $s=0$, its valuation is $3n$, giving the same contradiction. This proof applies also at $p=2$.
\end{proof}

\begin{proposition}\label{prop:cover}
The Jacobian of \eqref{eq:integral} is
\[
 E'_{AC}:\quad z^2=x^3+4(AC)^2x.
\]
The formulas
\[
 x=-At^2,\qquad z=AtY
\]
give a degree-two map from \eqref{eq:integral} to the congruent-number curve
\[
 E_{AC}:\quad z^2=x^3-(AC)^2x.
\]
Consequently a rational point with $tY\ne0$ is equivalent to a point on $E_{AC}(\Q)$ with abscissa in $-A\Q^{\times2}$ and nonzero ordinate.
\end{proposition}
\begin{proof}
Substitution gives
\[
 (AtY)^2=A^3t^2(C^2-t^4)=(-At^2)^3-(AC)^2(-At^2).
\]
Conversely an abscissa $x=-At^2\ne0$ determines $Y=z/(At)$. The quartic is the standard cover for the dual of the $2$-isogeny $E_{AC}\to E'_{AC}$ with kernel generated by $(0,0)$. Alternatively, the standard quartic Jacobian formula for $Y^2=a t^4+b$ gives $z^2=x^3-4abx$; here $a=-A,b=AC^2$, yielding the stated positive coefficient. The degree-two map identifies the quartic as a torsor under $E'_{AC}$.
\end{proof}

\begin{remark}\label{rem:gap}
Theorem~\ref{thm:local} places this cover in the relevant $2$-isogeny Selmer group. It does not make its class in $Section ha(E'_{AC})[2]$ vanish. Finiteness of that group is likewise insufficient. A sufficient additional hypothesis is $Section ha(E'_{AC})[2]=0$, together with positive Mordell--Weil rank to obtain infinitely many rational points and arbitrary finite guard avoidance. This is a precise extra hypothesis, not a consequence of ordinary abc and not an established uniform property of these twists.
\end{remark}

\begin{proposition}\label{prop:conditional}
Fix a prepared $d$. If $Section ha(E'_{AC})[2]=0$ and $E'_{AC}(\Q)$ has positive rank, then $C_d(\Q)$ is infinite. In particular, if $\mathcal A$ is a finite etale $\Q$-algebra containing a unit $d$ and a rational character sends $d$ to this prepared input, one may choose rational $s,y$ satisfying \eqref{eq:C} at that character and avoiding any specified finite set of guards whose norms are nonzero polynomials in $s$.
\end{proposition}
\begin{proof}
The genus-one curve has a rational divisor class of degree two, so its torsor class is killed by two. It is everywhere locally soluble by Theorem~\ref{thm:local}; the asserted vanishing of $Section ha[2]$ therefore gives a rational point. The curve then becomes isomorphic to its positive-rank Jacobian and has infinitely many rational points. The degree-two function $s$ has finite fibres, so there are infinitely many rational values of $s$. A finite collection of nonzero norm polynomials excludes only finitely many such values.
\end{proof}

For example the prepared input $d=112/169$ has the exact nonzero point
\[
 s=\frac{57}{104},\qquad y=\frac{4389}{35152}.
\]
This verifies one input, and does not establish uniform rational solubility.

\section{Relation to the two-coordinate block}
The intrinsic block is obtained from
\[
 p_0^2=d(1-e^2),\qquad
 q_0^2=d\bigl(s^4-e^2(1-d)^2\bigr).
\]
Setting $q_0=0$ and $e=\pm s^2/(1-d)$ yields \eqref{eq:C} with $y=(1-d)p_0$. The existing block's guard $q_0\ne0$ is not preserved by this specialization. Proposition~\ref{prop:sound} supplies a separate soundness proof for the replacement block; its quadratic root and guards must therefore be re-audited in the finite algebra. There is one free rational coordinate $s$ and one quadratic generator $y$. No uniform global realization, and hence no numerical improvement of the repository's bounds, is claimed.

\begin{thebibliography}{9}
\bibitem{XZ} M. Xiong and A. Zaharescu, \emph{Selmer groups and Tate--Shafarevich groups for the congruent number problem}, Comment. Math. Helv. 84 (2009), 21--56, Section 2.1. \url{https://doi.org/10.4171/CMH/151}
\bibitem{B} A. Burungale, \emph{A Diophantine definition of $\mathbb Z$ in $\mathbb Q$}, arXiv:2611.16130v1, Theorem B. \url{https://arxiv.org/html/2611.16130v1}
\end{thebibliography}
\end{document}
